% ---------------------------------------------------------------------------
% Phase 5.2 — F02 — Парная количественная оценка B-LoRA на SDXL.
%
% Аналог Table 1 оригинальной статьи B-LoRA: 50 пар (style, content) из
% манифеста DS8, SDXL backbone, DINO ViT-B/8 cosine similarity между выходом
% и (i) одиночным эталоном стиля, (ii) эталоном контента.
%
% Drop-in snippet для главы экспериментов. Требует booktabs + tabularx.
%
% Источники данных:
%   - Метрики:        output/results/f02_blora_sdxl_pairs.json
%   - Манифест пар:   experiments/data/b_lora_eval_pairs.json
%   - Visual sign-off:output/results/f02_blora_sdxl_pairs_visual.json
%   - Run timestamp:  20260515T152734
%
% Figure asset (требует ручной сборки, см. ниже):
%   figures/phase5_2_f02_grid.png
%   montage из 6 PNG в /tmp/review/f02_blora_sdxl_pairs/
% ---------------------------------------------------------------------------

\subsection{Парная количественная оценка B-LoRA на SDXL (протокол §5.1)}
\label{subsec:phase5_2_f02_quantitative}

Данный раздел воспроизводит методологию Table~1 оригинальной статьи B-LoRA в применении к собственному пайплайну SDXL\,+\,B-LoRA. Использован манифест \textbf{DS8}: восемь контентных объектов из репозитория DreamBooth (\textit{backpack}, \textit{bear}, \textit{bowl}, \textit{can}, \textit{cat}, \textit{clock}, \textit{dog}, \textit{vase}) и восемь стилевых референсов (по четыре полотна Клода Моне и Винсента ван Гога из WikiArt), из которых при $\mathtt{seed}=42$ зафиксировано \textbf{50} случайных пар $(\text{subject}, \text{style})$ без повторений. Для каждой пары генерируется одно изображение при склейке трёх LoRA-адаптеров (стилевой ствол ван Гога переиспользован из Phase~4.4, четыре стилевых ствола Моне и восемь контентных ствола обучены в текущем прогоне \texttt{run\_ts=20260515T152734}); метрики --- средние косинусные сходства DINO ViT-B/8 между выходом и одиночным стилевым эталоном (DINO-style) и между выходом и эталонным изображением контента (DINO-content).

В соответствии с зафиксированным в \texttt{experiments/datasets.md} \emph{reduced-protocol}-соглашением, абсолютные значения DINO-метрик \textbf{не сопоставимы напрямую с опубликованной B-LoRA Tab.~1}: оригинальная Table~1 опирается на полный пул $23\times25$ контент- и стиль-референсов, тогда как в настоящей работе используется уменьшенный пул $8\times8$ и подвыборка из 50 пар. Таким образом, отчёт ниже следует трактовать как методологическую реплику протокола (метрика, бэкбон, формула усреднения), а не как численное сравнение с литературой head-to-head.

\begin{table}[ht]
\centering
\caption{Парная оценка B-LoRA на SDXL по протоколу B-LoRA §5.1 (DS8, $\mathtt{seed}=42$, 50 пар). Метрики: DINO ViT-B/8 cosine между выходом и одиночным эталоном стиля / контента. STI~$=$~DINO-style~$-$~DINO-content.}
\label{tab:phase5_2_f02_dino}
\small
\begin{tabularx}{\textwidth}{l *{3}{Y}}
\toprule
\multicolumn{4}{c}{\textbf{Панель A. Сводка по художникам и в целом}} \\
\midrule
\textbf{Группа} & \textbf{DINO-style} $\uparrow$ & \textbf{DINO-content} $\uparrow$ & \textbf{STI} \\
\midrule
Van Gogh ($n=26$)     & \textbf{0{,}4215} & \textbf{0{,}2966} & $-0{,}1249$ \\
Claude Monet ($n=24$) & 0{,}3671          & 0{,}2573          & $-0{,}1098$ \\
\midrule
Всего ($n=50$)        & 0{,}3954          & 0{,}2777          & $\mathbf{-0{,}1177}$ \\
\bottomrule
\end{tabularx}

\vspace{0.6em}

\begin{tabularx}{\textwidth}{l *{3}{Y}}
\toprule
\multicolumn{4}{c}{\textbf{Панель B. Разбиение по стилевому LoRA} ($n$ --- число пар, в которых данный стиль участвует)} \\
\midrule
\textbf{style\_id} & $\boldsymbol{n}$ & \textbf{DINO-style} $\uparrow$ & \textbf{DINO-content} $\uparrow$ \\
\midrule
wikimo\_01 & 8 & 0{,}3222 & 0{,}3006 \\
wikimo\_02 & 5 & \textbf{0{,}5179}\textsuperscript{\dag} & 0{,}2071 \\
wikimo\_03 & 6 & 0{,}3308 & 0{,}2674 \\
wikimo\_04 & 5 & 0{,}3317 & 0{,}2261 \\
\midrule
wikivg\_01 & 6 & 0{,}4401 & 0{,}2705 \\
wikivg\_02 & 6 & \textbf{0{,}5065} & 0{,}2891 \\
wikivg\_03 & 7 & 0{,}3264 & 0{,}2968 \\
wikivg\_04 & 7 & 0{,}4278 & 0{,}3253 \\
\bottomrule
\end{tabularx}

\vspace{0.6em}

\begin{tabularx}{\textwidth}{l *{3}{Y}}
\toprule
\multicolumn{4}{c}{\textbf{Панель C. Разбиение по контентному субъекту}} \\
\midrule
\textbf{content\_subject\_id} & $\boldsymbol{n}$ & \textbf{DINO-style} $\uparrow$ & \textbf{DINO-content} $\uparrow$ \\
\midrule
backpack & 7 & 0{,}3368 & 0{,}3171 \\
bear     & 5 & 0{,}3332 & 0{,}1583 \\
bowl     & 8 & 0{,}4521 & 0{,}1547 \\
can      & 6 & 0{,}4406 & 0{,}3900 \\
cat      & 5 & 0{,}4561 & \textbf{0{,}4584} \\
clock    & 6 & 0{,}3132 & 0{,}2860\textsuperscript{\dag} \\
dog      & 5 & \underline{0{,}2258} & 0{,}2528 \\
vase     & 8 & \textbf{0{,}5245} & 0{,}2532 \\
\bottomrule
\end{tabularx}

\vspace{0.4em}
\footnotesize
\textsuperscript{\dag} Ячейка с \texttt{pair\_id}~$=$~4 (\texttt{wikimo\_02}~$\times$~\texttt{clock}, DINO-content~$=-0{,}013$) --- failure-mode \emph{style-LoRA over-dominance}: предмет (часы) на выходе отсутствует, кадр редуцирован к стилевому прототипу <<Кувшинки>> Моне. Ячейка учтена в средних значениях по строкам \texttt{wikimo\_02} и \texttt{clock}, отдельно проанализирована в качественном разделе (см.~\autoref{fig:phase5_2_f02_grid}). Полужирным выделены максимумы по соответствующей колонке внутри панели, подчёркиванием --- глобальный минимум DINO-style по субъектам. \textbf{Важно:} абсолютные значения DINO в данной таблице не сопоставимы с опубликованной B-LoRA Tab.~1 (полный пул $23\times25$); настоящий протокол --- reduced ($8\times8$, 50 пар), см.~\texttt{experiments/datasets.md}.
\end{table}

\normalsize

Введём интегральный показатель баланса осей по аналогии с Phase~6:
\begin{equation}
\mathrm{STI}_{\text{SDXL,F02}}
   = \overline{\mathrm{DINO\text{-}style}} - \overline{\mathrm{DINO\text{-}content}}
   = 0{,}3954 - 0{,}2777
   = -0{,}1177.
\label{eq:phase5_2_sti}
\end{equation}
Отрицательное значение $\mathrm{STI}\approx -0{,}118$ свидетельствует о слабом смещении системы в сторону сохранения содержания: контент <<выигрывает>> у стиля приблизительно на $0{,}12$ единицы косинусного сходства. По модулю это значение существенно меньше, чем у конфигурации FLUX из Phase~6 ($\mathrm{STI}_{\mathrm{FLUX,M01b3}}=+0{,}463$, см.~\eqref{eq:phase6_sti}), причём знак противоположен: SDXL-конфигурация одновременно поднимает стилевую ось и снижает контентную, приближая баланс к нейтральному. Это согласуется с тезисом литературы о том, что иерархическая структура U-Net в SDXL допускает более чистую изоляцию стилевых блоков (\texttt{up\_blocks.0.attentions.1}), чем глобальные двойные стримы FLUX.

Разбиение по художнику фиксирует устойчивый разрыв в пользу Van Gogh: $\overline{\mathrm{DINO\text{-}style}}_{\mathrm{vg}} = 0{,}4215$ против $0{,}3671$ для Monet, и $\overline{\mathrm{DINO\text{-}content}}_{\mathrm{vg}} = 0{,}2966$ против $0{,}2573$. Превосходство по обеим осям одновременно делает невозможным объяснение через trade-off; более правдоподобной является гипотеза о различной <<обучаемости>> стилей low-rank адаптацией: полотна Van Gogh содержат высокочастотные мазки выраженной геометрической формы и насыщенную палитру, тогда как стиль Monet опирается на тонкие тональные градации и пастельные переходы, хуже кодируемые в LoRA ранга $r=16$. Дополнительно играет роль происхождение адаптеров: четыре Van Gogh LoRA переиспользованы из Phase~4.4 с длительным циклом отбора и валидации, тогда как четыре Monet LoRA обучены в текущем прогоне в режиме паритета (500 шагов, $r=16$) без дополнительной настройки.

На уровне отдельных стилей наибольший стилевой отклик демонстрируют \texttt{wikimo\_02} (<<Мадам Моне с ребёнком>>, $\overline{\mathrm{DINO\text{-}style}}=0{,}5179$) и \texttt{wikivg\_02} (<<Автопортрет с перевязанным ухом>>, $0{,}5065$), наименьший --- \texttt{wikimo\_01} (<<Стог сена в Живерни>>, $0{,}3222$) и \texttt{wikivg\_03} (<<Подсолнухи>>, $0{,}3264$). Аномально высокий показатель \texttt{wikimo\_02} частично объясняется зафиксированным failure-кейсом (\texttt{pair\_id}$=$4), завышающим столбец за счёт почти полного исчезновения контента. Среди контентных субъектов наиболее устойчивыми носителями стиля оказываются \textit{vase} ($0{,}5245$), \textit{cat} ($0{,}4561$) и \textit{bowl} ($0{,}4521$) --- объекты с простой осевой симметрией или богатой собственной семантикой, оставляющие достаточно пиксельного <<пространства>> для стилевого отпечатка. Худший стилевой перенос наблюдается у \textit{dog} ($\overline{\mathrm{DINO\text{-}style}}=0{,}2258$), что коррелирует с категорией одушевлённости: предобученный на ImageNet U-Net удерживает сильный семантический prior на собаке как фотореалистичной форме, и низкоранговая LoRA-инъекция стиля не способна его полностью переписать.

Отдельного обсуждения заслуживает пара \texttt{pair\_id}$=$4 (\texttt{wikimo\_02}~$\times$~\texttt{clock}) с DINO-content~$=-0{,}013$ --- единственный полный content-collapse в выборке. Визуальный sign-off подтверждает, что вместо часов в кадре генерируется типовая монетовская композиция кувшинок на пруду; косинусное сходство с эталонным изображением часов обнуляется и формально уходит в отрицательную область шума. Данный кейс сохранён в выборке без купирования и трактуется как иллюстрация режима <<style-LoRA over-dominance>>, при котором стилевой адаптер при слиянии получает эффективный масштаб, превышающий контентный, и подавляет семантический сигнал. Подробный визуальный разбор приведён в качественном разделе главы.

Сопоставление с Phase~5.1 фиксирует ожидаемое снижение grand-mean при переходе от одностилевой к парной оценке: для winner-конфига Phase~5.1 \texttt{F01-4} (img4) на той же SDXL-архитектуре было получено DINO-style $=0{,}4208$, тогда как F02 на DS8 даёт DINO-style $=0{,}3954$ --- падение примерно на $-6\%$. Объяснение методологическое: DS8 включает восемь различных контентных категорий, в том числе сложные для текстурного переноса (\textit{clock}, \textit{dog}), и шесть случайных стиль$\times$контент комбинаций на каждый стилевой LoRA, тогда как Phase~5.1 измерял стиль на одной курированной серии. Таким образом, наблюдаемое снижение grand-mean относительно Phase~5.1 является ожидаемым следствием расширения протокола, а не регрессией метода; per-artist срез Van Gogh ($\overline{\mathrm{DINO\text{-}style}}=0{,}4215$) фактически воспроизводит число Phase~5.1 в пределах третьего знака.

\subsection{Качественная иллюстрация поведения B-LoRA на SDXL}
\label{subsec:phase5_2_f02_qualitative}

\begin{figure}[ht]
    \centering
    \includegraphics[width=\textwidth]{figures/phase5_2_f02_grid.png}
    \caption{Качественная иллюстрация поведения B-LoRA на SDXL в рамках протокола DS8 (Phase 5.2 / F02, 50 пар, run\_ts 20260515T152734): сетка $2 \times 3$, по которой отобраны характерные ячейки из общей выборки. Верхняя строка --- успешные пары с одновременным переносом стиля и сохранением содержания: \textit{wikivg\_02} (Ван Гог) $\times$ \textit{bowl} (DINO-style $= 0{,}583$, DINO-content $= 0{,}115$), \textit{wikivg\_02} $\times$ \textit{can} ($0{,}523$ и $0{,}373$ соответственно), \textit{wikivg\_02} $\times$ \textit{cat} ($0{,}522$ и $0{,}498$). Нижняя строка --- характерные ограничения метода: \textit{wikimo\_03} (Моне) $\times$ \textit{backpack} (DINO-style $= 0{,}296$, DINO-content $= 0{,}313$) с уравновешенным, но визуально приглушённым стилем; \textit{wikimo\_03} $\times$ \textit{clock} ($0{,}172$ и $0{,}356$), где стилевая текстура локализована внутри контура предмета вместо глобального наложения; \textit{wikimo\_02} $\times$ \textit{clock} (DINO-style $= 0{,}489$, DINO-content $= -0{,}013$) --- канонический режим отказа <<content-collapse>>, при котором стилевой адаптер полностью подавляет геометрию объекта и порождает изображение, близкое к фрагменту обучающего полотна. Числовые подписи воспроизводят покомпонентные значения DINO ViT-B/8, агрегированные в \autoref{tab:phase5_2_f02_dino} (большое среднее DINO-style $= 0{,}3954$, DINO-content $= 0{,}2777$, $\mathrm{STI} = -0{,}118$).}
    \label{fig:phase5_2_f02_grid}
\end{figure}

Как видно из \autoref{fig:phase5_2_f02_grid}, перенос стиля на SDXL в рамках протокола DS8 значимо асимметричен по выбору стилевого источника: ячейки с \textit{wikivg\_02} (верхняя строка) демонстрируют устойчивое наложение характерной мазочной структуры Ван Гога при сохранении распознаваемой геометрии контента, тогда как пары с \textit{wikimo\_*} (нижняя строка) систематически вырождаются либо в приглушённую цветовую перекраску без выраженного мазка, либо в режим <<style-on-canvas>>, когда модель воспроизводит фрагмент обучающего пейзажа Моне поверх запрошенного объекта. Наиболее наглядный случай этого режима представлен в нижней строке, правая ячейка (\textit{wikimo\_02} $\times$ \textit{clock}): объект-часы вытесняется садово-водной композицией стиля, что количественно фиксируется аномальным значением DINO-content $= -0{,}013$, отмеченным в~\autoref{tab:phase5_2_f02_dino} как глобальный минимум по протоколу и интерпретируемым как канонический failure-mode <<content-collapse>> при доминировании style-LoRA. Совокупно эти наблюдения согласуются с интегральным показателем $\mathrm{STI} = -0{,}118$, то есть с инверсией баланса осей по сравнению с FLUX-конфигурацией Phase 6, и обосновывают необходимость отдельной настройки масштабов стилевого и контентного каналов на SDXL.
